% Figura: método de los trabajos virtuales. Desplazamiento infinitesimal de un
% conductor (a) aislado (cargas constantes) y (b) conectado a baterías (potenciales constantes).
\begin{figure}[t]
\centering
\tikzset{
  cond/.style={fill=conductor, draw=texto, line width=0.6pt},
  condv/.style={draw=texto, dashed, line width=0.5pt, dash pattern=on 2.5pt off 1.8pt},
  bateria/.pic={
    \draw[texto, line width=0.6pt] (-0.35,0) -- (-0.08,0) (0.08,0) -- (0.35,0);
    \draw[texto, line width=1.1pt] (-0.08,-0.22) -- (-0.08,0.22);
    \draw[texto, line width=0.6pt] (0.08,-0.13) -- (0.08,0.13);}
}
\begin{subfigure}[b]{0.47\textwidth}\centering
\begin{tikzpicture}[scale=0.88]
  \draw[cond] plot[smooth cycle, tension=0.9] coordinates {(-2.3,0.1) (-1.9,0.75) (-1.2,0.6) (-1.0,-0.1) (-1.6,-0.6) (-2.2,-0.45)};
  \node[font=\small, text=black] at (-1.65,0.05) {$Q_1$};
  \draw[cond] plot[smooth cycle, tension=0.9] coordinates {(0.6,0.25) (1.0,0.85) (1.6,0.7) (1.75,0.0) (1.3,-0.55) (0.75,-0.4)};
  \node[etiq, anchor=south] at (1.2,0.9) {$Q_2$};
  \draw[condv] plot[smooth cycle, tension=0.9] coordinates {(1.05,0.25) (1.45,0.85) (2.05,0.7) (2.2,0.0) (1.75,-0.55) (1.2,-0.4)};
  \draw[vecP] (1.75,0.0) -- (2.2,0.0);
  \node[etiq, text=campoP, anchor=west] at (2.25,0.0) {$\dl$};
  \node[etiqt, anchor=north] at (-0.3,-1.0) {cargas aisladas: $\dd Q_i=0$};
  \node[etiq, anchor=north] at (-0.3,-1.45) {$\dd W_{\mathrm{mec}}+\dd U_e=0$};
\end{tikzpicture}
\caption{Sistema aislado}
\end{subfigure}\hfill
\begin{subfigure}[b]{0.47\textwidth}\centering
\begin{tikzpicture}[scale=0.88]
  \draw[cond] plot[smooth cycle, tension=0.9] coordinates {(-2.3,0.1) (-1.9,0.75) (-1.2,0.6) (-1.0,-0.1) (-1.6,-0.6) (-2.2,-0.45)};
  \node[font=\small, text=black] at (-1.65,0.05) {$\phi_1$};
  \draw[cond] plot[smooth cycle, tension=0.9] coordinates {(0.6,0.25) (1.0,0.85) (1.6,0.7) (1.75,0.0) (1.3,-0.55) (0.75,-0.4)};
  \node[etiq, anchor=south] at (1.2,0.9) {$\phi_2$};
  \draw[condv] plot[smooth cycle, tension=0.9] coordinates {(1.05,0.25) (1.45,0.85) (2.05,0.7) (2.2,0.0) (1.75,-0.55) (1.2,-0.4)};
  \draw[vecP] (1.75,0.0) -- (2.2,0.0);
  % cables y baterías
  \draw[texto, line width=0.55pt] (-1.65,-0.55) -- (-1.65,-1.0) -- (-0.75,-1.0);
  \pic at (-0.4,-1.0) {bateria};
  \draw[texto, line width=0.55pt] (-0.05,-1.0) -- (0.25,-1.0);
  \draw[texto, line width=0.55pt] (1.15,-0.5) -- (1.15,-1.0) -- (0.55,-1.0);
  \pic at (0.4,-1.0) {bateria};
  \node[etiq, text=campoP, anchor=west] at (2.25,0.0) {$\dl$};
  \node[etiqt, anchor=north] at (-0.1,-1.35) {potenciales fijos: $\dd\phi_i=0$};
  \node[etiq, anchor=north] at (-0.1,-1.8) {$\dd W_{\mathrm{mec}}+\dd U_e=\dd W_{\mathrm{bat}}$};
\end{tikzpicture}
\caption{Sistema a potencial constante}
\end{subfigure}
\caption{Método de los trabajos virtuales. Se imagina un desplazamiento infinitesimal $\dl$ de uno de los conductores
(trazos) y se hace el balance de energía. (a) Si el sistema está aislado, las cargas no cambian y el trabajo mecánico
$\dd W_{\mathrm{mec}}=\vF_e\cdot\dl$ sale de la energía electrostática. (b) Si los conductores están unidos a baterías, sus potenciales
no cambian, pero las baterías deben suministrar las cargas $\dd Q_i$, con un trabajo $\dd W_{\mathrm{bat}}=\sum_i\phi_i\,\dd Q_i$.}
\label{fig:desplazamiento}
\end{figure}
